\section{Proof of the main theorem}\label{sec:main-proof}

The upcoming Parts~\ref{part:types}, \ref{part:definable-construction-sequences}, and~\ref{part:mc} prove three independent results.
This section combines them into the main result, \Cref{thm:main}. It does so via two theorems that are connected by the following type-based width notion.

% lax begin Lax185574.Width
\begin{restatable}{definition}{defrqwidth}\label{res:width}
For $r,q\in\N$,
the \emph{radius-$r$ depth-$q$ width}
of a construction sequence $G_1,\ldots,G_\lambda$
is defined as:
   $$\max_{t\in[\lambda-1]}\max_{v\in V(G)}|\Types^{G_t}_{q,q}({\Ball^{G_{t+1}}_r(v)})|.$$
  \end{restatable}
% lax end

The offset between $G_t$ and $G_{t+1}$ mirrors the offset between $\PP_{i-1}$ and $R_i$ in the width definition of chains and of merge sequences (see \Cref{def:chain-width,def:merge-sequence}).

% lax begin Lax185574.FewTypes
\begin{theorem}\label{thm:skel-few-types}
  Let $\CC$ be a monadically dependent graph class, $r,q\in\N$, and let $\eps>0$ be a rational.
  There is a constant $\alpha\ge1$ depending only on $\CC$ and $q$ such that every $n$-vertex graph $G\in\CC$
  has a construction sequence of \text{radius-$r$} depth-$q$ width $O_{\CC,r,q,\eps}(n^{\alpha\eps})$ and length $O_{\CC,r,q,\eps}(n^{1+\alpha\eps})$.
  Moreover, there is an algorithm (independent of~$\CC$) which, given $r,q\in\N$, a rational $\eps>0$, and an $n$-vertex graph $G\in\CC$,
  computes such a construction sequence in time $\widehat O_{\CC,r,q,\eps}(n^3)$.
\end{theorem}
% lax end
\noindent
It is important that $\alpha$ does not depend on $\eps$, as this allows us to make the exponent $\alpha\eps$ of the width as small as we like.
The upcoming \Cref{sec:skel-few-types} assembles \Cref{thm:skel-few-types} from the key results of
\Cref{part:types,part:definable-construction-sequences}.

\newcommand{\mcdepth}{q}
% lax begin Lax185574.ConstructionModelChecking
\begin{restatable}{theorem}{mcthm}\label{thm:mc}
  There is an algorithm which, given
  a sentence $\phi$ and an $n$-vertex graph~$G$ together with its construction sequence of length $\lambda$, determines whether $G\models\phi$ in time
  $\lambda\cdot n^2\cdot\poly_{\phi}(K)$,
  where $K$ is the radius-$r$ depth-$\mcdepth$ width of the construction sequence,
  for some $r$ and $\mcdepth$ which are computable from $\phi$.
\end{restatable}

% lax end

\begin{proof}
  See \Cref{part:mc}.
\end{proof}

From these two theorems, we can derive
\Cref{thm:main} without much effort.
The only subtlety is that the algorithm does not know $\CC$, so it cannot choose $\eps$ itself.

\mainthm*

% lax begin Lax185574Proofs.ModelChecking.exists_program_modelChecking
\begin{proof}[Proof of \Cref{thm:main}]
  Let $\phi$ be the input sentence, and let $r,q$ be the numbers from \Cref{thm:mc}.
  For $j\ge1$, let $\mathcal A_j$ be the algorithm that computes a construction sequence of $G$ by \Cref{thm:skel-few-types} with $r,q$, and $\eps\coloneqq1/j$,
  and then decides whether $G\models\phi$ by \Cref{thm:mc}.

  Fix a monadically dependent graph class $\CC$, and let $\alpha\ge1$ be the constant from \Cref{thm:skel-few-types} for $\CC$ and $q$.
  Let $c\ge1$ be such that the polynomial in \Cref{thm:mc} is $O_\phi(K^c)$.
  On $n$-vertex graphs $G\in\CC$, the sequence computed by $\mathcal A_j$ has radius-$r$ depth-$q$ width $K\le O_{\CC,\phi,j}(n^{\alpha/j})$ and length $\lambda\le O_{\CC,\phi,j}(n^{1+\alpha/j})$.
  So $\mathcal A_j$ runs in time
  \[
    \widehat O_{\CC,\phi,j}(n^3)+O_\phi(\lambda\cdot n^2\cdot K^c)
    \le O_{\CC,\phi,j}(n^{3+\beta/j}),
  \]
  where $\beta\coloneqq(c+1)\alpha$.
  Since $\beta/j>0$, the subpolynomial factor hidden by $\widehat O$ is absorbed by $n^{\beta/j}$.
  The number $\beta$ depends on $\CC$ and $\phi$, but not on~$j$.

  The algorithm does not know $\CC$, so it cannot choose $j$.
  Instead, it runs all algorithms $\mathcal A_j$ in parallel, giving $\mathcal A_j$ a $2^{-j}$ share of the computation, and returns the first answer.
  Its answer is correct, as every $\mathcal A_j$ correctly solves the model checking problem on every input graph.

  For the running time, let $\eps>0$ and $j\coloneqq\lceil2\beta/\eps\rceil$.
  If $\mathcal A_j$ halts within time $T$, then it completes within the first $2^jT$ simulated steps, so the algorithm halts within time
  \[
    O\bigl(2^jT\log(2^jT+2)\bigr)=O_j\bigl(T\log(T+2)\bigr).
  \]
  Hence, on $n$-vertex graphs in $\CC$, the algorithm runs in time
  \[
    O_{\CC,\phi,\eps}\bigl(n^{3+\beta/j}\log(n+2)\bigr)
    \le O_{\CC,\phi,\eps}(n^{3+\eps}),
  \]
  since $\beta/j\le\eps/2$ and $\log(n+2)=O_\eps(n^{\eps/2})$.
\end{proof}
% lax end

\subsection{Proving \texorpdfstring{\Cref{thm:skel-few-types} from \Cref{part:types,part:definable-construction-sequences}}{Theorem \getrefnumber{thm:skel-few-types} from Parts I and II}}\label{sec:skel-few-types}

We first state the key results of the two parts, and derive \Cref{thm:skel-few-types} afterwards.

\subsubsection{The key result of \texorpdfstring{\Cref{part:types}}{Part I}}\label{sec:skel-part-types}

The key ingredient towards proving \Cref{thm:skel-few-types} is the following type bound derived in \Cref{part:types}.

% lax begin Lax185574.TypeBound
\newif\iftypeboundgeneral
\begin{restatable}{theorem}{thmtypes}\label{thm:type-bound}

Let $\CC$ be a monadically dependent graph class and $k,q,c\in\N$.
For every \iftypeboundgeneral(possibly infinite) \fi graph $G \in \CC$, expansion $G^+$ of $G$ by $c$ unary predicates, and \iftypeboundgeneral finite \fi subset $U \subseteq V(G)$,
\[
    \left|\Types^{G^+}_{k,q}(V(G)/U)\right| \leq O_{\CC,q,k,c}(1)\cdot \poly_{\CC,q,k}(|U|).
\]

\end{restatable}
% lax end

The number $c$ of unary predicates affects only the constant factor, not the degree of the polynomial.
This matters, as our construction sequences use $O(r/\eps)$ unary predicates, while the constant $\alpha$ in \Cref{thm:skel-few-types} must not depend on $\eps$.

\subsubsection{The key result of \texorpdfstring{\Cref{part:definable-construction-sequences}}{Part II}}\label{sec:skel-part-defseq}

The type bound counts types in $G$, while the width (\Cref{res:width}) counts types in the partial graphs $G_t$.
To bridge this gap, \Cref{part:definable-construction-sequences} produces a construction sequence in which local balls of $G_t$ are defined by formulas in $G$.

% lax begin Lax185574.LocalDefinability
\begin{restatable}{definition}{defpartialdef}\label{def:partial-definable}
Let $G$ be a graph with vertex set $V$ and let $H$ be a partial graph with vertex set $W\subseteq V$.
For $q,c,w\in \N$,
we say that $H$ is \emph{definable} in $G$  by formulas of quantifier rank $q$ using $c$ unary predicates and $w$ parameters, if
there is an expansion $G^+$ of $G$ by $c$ unary predicates
and a set $S\subseteq V(G)$ of parameters with $|S|\le w$,
such that for every relation symbol $Q$ of $H$ (when viewed as a relational structure) of arity $k$ there is some formula $\phi_Q(x_1,\ldots,x_k)$ of quantifier rank $q$ with parameters from $S$
such that
$$G^+\models\phi_Q(\bar v)\quad\iff\quad \bar v\in W^k\text{ and }H\models Q(\bar v)\qquad\quad\text{for all $\bar v\in V^k$.}$$
\end{restatable}

The vertex set $W$ of $H$ is the union of its parts, so it is defined by the disjunction of the formulas defining the parts.

\begin{restatable}{definition}{defskellocdef}\label{def:skel-loc-def}
  Let $G$ be a graph and $r,c,s,q_0\in\N$.
  A construction sequence $G_1,\ldots,G_\lambda$ of $G$ is \emph{radius-$r$ definable} with $c$ unary predicates, $s$ parameters, and quantifier rank $q_0$
  if for all $t\in[\lambda-1]$ and $v\in V(G)$, the induced partial graph
  \[
      G_t[\Ball^{G_{t+1}}_r(v)]
  \]
  is definable in $G$ by formulas of quantifier rank $q_0$ using $c$ unary predicates and $s$ parameters.
\end{restatable}

\noindent Note that the expansion and the parameters may depend on $t$ and $v$.
% lax end

\medskip

The key result of \Cref{part:definable-construction-sequences} is the following theorem.
% lax begin Lax185574.DefinableSequences
\begin{restatable}{theorem}{thmskeldefseq}\label{thm:skel-def-seq}
  Let $r\ge 1$ be an integer, $\eps>0$ be a rational, and $\CC$ be a monadically dependent graph class. There is a constant $\delta\ge1$ depending only on $\CC$ such that every $n$-vertex graph $G\in\CC$
  has a construction sequence of $G$ which is radius-$r$ definable with $O(r/\eps)$ unary predicates, $O_{\CC,r,\eps}(n^{\delta\eps})$ parameters, and quantifier rank $7$, and has length $O_{\CC,r,\eps}(n^{1+\delta\eps})$. Moreover, there is an algorithm (independent of $\CC$)
  which, given $r$, $\eps$, and $G$, computes such a sequence, and runs in time $\widehat O_{\CC,r,\eps}(n^3)$ whenever $G\in\CC$.
\end{restatable}
% lax end

For nowhere dense classes, the following stronger bounds follow easily from known results. The proof is given in \Cref{sec:nd}.

\begin{restatable}{theorem}{thmnddefseq}\label{thm:nd-def-seq}
  Let $\CC$ be a nowhere dense graph class, $r\in\N$, and $\eps>0$.
  Every $n$-vertex graph $G\in\CC$ has a radius-$r$ definable construction sequence with one unary predicate, $O_{\CC,r,\eps}(n^\eps)$ parameters, quantifier rank $0$, and length $O_{\CC,r,\eps}(n^{1+\eps})$.
  Such a sequence can be computed in time $O_{\CC,r,\eps}(n^{1+\eps})$.
\end{restatable}

In particular, \Cref{thm:nd-def-seq}, together with the results of \Cref{part:types,part:mc}, gives an alternative proof of the result of Grohe, Kreutzer, and Siebertz~\cite{gks} that first-order model checking is fixed-parameter tractable on nowhere dense classes.
(Alternatively, in nowhere dense classes one can replace the type bound from \Cref{part:types}  by a stronger bound deducible from \cite{number-of-types}).

\subsubsection{Assembling the parts}\label{sec:skel-assembly}

We use the following consequence of Gaifman locality, proved in \Cref{app:gaifman}.

% lax begin Lax185574.LocalToGlobal
\begin{restatable}[Corollary of Gaifman locality]{lemma}{lemLocalToGlobal}\label{lem:local-to-global}
For every $k,q,\alpha \in\N$, there are numbers $h=h(q)$ and $q'=q'(k,q,\alpha)$ computable from $k,q,\alpha$,
 with the following property.
Fix a finite relational signature $\sigma$ of arity at most $\alpha$ and a $\sigma$-structure $A$.
Then, for all $\bar a\in V(A)^k$
and $B\supseteq \Ball_h^A(\bar a)$,
$\tp_q^A(\bar a)$ depends only on
 $\tp_{q'}^{A[B]}(\bar a)$.
\end{restatable}
% lax end

% lax begin Lax185574.DefinableToTypes
\begin{lemma}\label{lem:skel-def-to-types}
  For every $q\in\N$, there is $h\in\N$, computable from $q$, such that the following holds for every monadically dependent graph class $\CC$ and all $r,c,s,q_0\in\N$.
  If a construction sequence of a graph $G\in\CC$ is radius-$(r+h)$ definable with $c$ unary predicates, $s$ parameters, and rank $q_0$,
  then its radius-$r$ depth-$q$ width is at most $O_{\CC,q,q_0,c}(1)\cdot\poly_{\CC,q,q_0}(s)$.
\end{lemma}
% lax end
% lax begin Lax185574Proofs.DefinableToTypes.width_le_of_isRadiusDefinable
\begin{proof}
  We will apply \Cref{lem:local-to-global} to partial graphs from a construction sequence of $G$.
  These partial graphs will have a finite signature $\sigma$ with up to $|V(G)|$ unary predicates marking the partition, but the arity is fixed to $\alpha=2$.
  Hence, we can choose $h\coloneqq h(q)$ and $q'\coloneqq q'(q,q,2)$ from \Cref{lem:local-to-global}.
  Let $G_1,\ldots,G_\lambda$ be a construction sequence of $G\in\CC$ which is radius-$(r+h)$ definable with $c$ unary predicates, $s$ parameters, and rank $q_0$.
  Fix $t\in[\lambda-1]$ and $v\in V(G)$, and let $B\coloneqq\Ball^{G_{t+1}}_r(v)$ and $W\coloneqq\Ball^{G_{t+1}}_{r+h}(v)$.
  By \Cref{res:width}, we have to bound $\left|\Types^{G_t}_{q,q}(B)\right|$.

  First, we move from $G_t$ to the partial graph $G_t[W]$.
  As $G_{t+1}$ has all resolved pairs of $G_t$, every tuple $\bar u\in B^q$ satisfies $\Ball^{G_t}_h(\bar u)\subseteq W$.
  So by \Cref{lem:local-to-global}, $\tp^{G_t}_q(\bar u)$ is determined by $\tp^{G_t[W]}_{q'}(\bar u)$, and hence
  \[
    \left|\Types^{G_t}_{q,q}(B)\right|
    \le\left|\Types^{G_t[W]}_{q,q'}(W)\right|.
  \]

  Next, we move from $G_t[W]$ to $G$.
  By assumption, $G_t[W]$ is definable in an expansion $G^+$ of $G$ by $c$ unary predicates, by formulas of quantifier rank $q_0$ with parameters from a set $S$ of at most $s$ vertices.
  Replace every atom of a formula over $G_t[W]$ by its defining formula, and restrict every quantifier to the vertex set of $G_t[W]$.
  This turns a formula of quantifier rank $q'$ over $G_t[W]$ into a formula of quantifier rank at most $q'+q_0$ over $G^+$ with parameters from $S$,
  which holds for the same tuples from $W$.
  So for every $\bar u\in W^q$, $\tp^{G_t[W]}_{q'}(\bar u)$ is determined by $\tp^{G^+}_{q'+q_0}(\bar u/S)$.

  Together with \Cref{thm:type-bound}, we obtain
  \[
    \left|\Types^{G_t}_{q,q}(B)\right|
    \le\left|\Types^{G_t[W]}_{q,q'}(W)\right|
    \le\left|\Types^{G^+}_{q,q'+q_0}(V(G)/S)\right|
    \le O_{\CC,q,q'+q_0,c}(1)\cdot\poly_{\CC,q,q'+q_0}(s).
  \]
  As $q'$ depends only on $q$, this is the claimed bound $O_{\CC,q,q_0,c}(1)\cdot\poly_{\CC,q,q_0}(s)$.
\end{proof}
% lax end

% lax begin Lax185574Proofs.FewTypes.exists_program_fewTypes
\begin{proof}[Proof of \Cref{thm:skel-few-types}]
  Let $\CC$ be a monadically dependent graph class containing $G$.
  Let $h$ be the number given by \Cref{lem:skel-def-to-types} for depth $q$.
  We run the algorithm of \Cref{thm:skel-def-seq} with radius $r+h$ and $\eps$, which takes time $\widehat O_{\CC,r,q,\eps}(n^3)$.

  By \Cref{thm:skel-def-seq}, there is a constant $\delta\ge1$ depending only on $\CC$ such that the sequence is radius-$(r+h)$ definable with $c\le O((r+h)/\eps)\le O_{r,q,\eps}(1)$ unary predicates, $s\le O_{\CC,r,q,\eps}(n^{\delta\eps})$ parameters, and rank $q_0\coloneqq7$,
  and it has length $O_{\CC,r,q,\eps}(n^{1+\delta\eps})$.
  By \Cref{lem:skel-def-to-types}, there is an integer $D\ge1$ such that its radius-$r$ depth-$q$ width is at most
  \[
    O_{\CC,q,q_0,c}(1)\cdot\poly_{\CC,q,q_0}(s)\le O_{\CC,r,q,\eps}(n^{D\delta\eps}),
  \]
  where we used $q_0=7$ and $c\le O_{r,q,\eps}(1)$. As $D$ depends only on $\CC$ and $q$, we can choose $\alpha\coloneqq D\delta$ to satisfy the statement.
  Since $D\ge1$ and $\delta\ge1$, we have $\alpha\ge \delta\ge1$, so this choice also bounds the sequence length by $O_{\CC,r,q,\eps}(n^{1+\alpha\eps})$.
\end{proof}
% lax end
